\subsubsection{Abraham--Lorentz--Dirac 方程的可解应用、无 runaway 条件与降阶}

前面已经给出点电荷的完整协变方程。为看清这个三阶世界线方程的解空间，先不作低速近似，而考虑严格的一维运动：用快度 $\psi(\tau)$ 参数化四速度，
\begin{equation*}
 u^\mu=c\bigl(\cosh\psi,\sinh\psi,0,0\bigr),
 \qquad
 v=c\tanh\psi,
 \qquad
 p=mc\sinh\psi.
\end{equation*}
对固有时逐次求导，得到
\begin{align*}
 a^\mu
 &=c\dot\psi\bigl(\sinh\psi,\cosh\psi,0,0\bigr),\\
 \dot a^\mu
 &=c\ddot\psi\bigl(\sinh\psi,\cosh\psi,0,0\bigr)
 +c\dot\psi^2\bigl(\cosh\psi,\sinh\psi,0,0\bigr),\\
 a^2&=-c^2\dot\psi^2.
\end{align*}
其中第三式是 $a^\mu$ 自缩并的结果，用到 $\sinh^2\psi-\cosh^2\psi=-1$。把最后一式代入 ALD 组合，含 $\dot\psi^2$ 的两项逐分量抵消：
\begin{align*}
 \dot a^\mu+{a^2\over c^2}u^\mu
 &={}c\ddot\psi(\sinh\psi,\cosh\psi,0,0)
 +c\dot\psi^2(\cosh\psi,\sinh\psi,0,0)
 -c\dot\psi^2(\cosh\psi,\sinh\psi,0,0)\\
 &=c\ddot\psi(\sinh\psi,\cosh\psi,0,0).
\end{align*}
令 $F_\parallel(\tau)$ 表示实验室中的一维普通力，即 $F_\parallel=\d p/\d t$。由 $v=c\tanh\psi$ 解出 $\gamma=(1-\tanh^2\psi)^{-1/2}=\cosh\psi$、$\gamma v/c=\sinh\psi$，相应四力为
\begin{align*}
 f_{\mathrm{ext}}^\mu
 &=\gamma\left({F_\parallel v\over c},F_\parallel,0,0\right)\\
 &=F_\parallel(\sinh\psi,\cosh\psi,0,0).
\end{align*}
以上三个式子中，惯性项、外力项和自力项都平行于同一个单位类空方向 $(\sinh\psi,\cosh\psi,0,0)$。快度参数化的威力正在于此：三条平行方向使四维向量方程退化为单个标量方程，因此四维 ALD 方程严格化为
\begin{equation*}
 mc\left(\dot\psi-\tau_{\mathrm{rr}}\ddot\psi\right)
 =F_\parallel(\tau),
 \qquad
 \tau_{\mathrm{rr}}={q^2\over6\pi\varepsilon_0mc^3}.
\end{equation*}
这一化简没有使用 $v/c\ll1$：把 $a^\mu=c\dot\psi(\sinh\psi,\cosh\psi,0,0)$、$\dot a^\mu+a^2u^\mu/c^2=c\ddot\psi(\sinh\psi,\cosh\psi,0,0)$ 与四力代入 $ma^\mu=f_{\mathrm{ext}}^\mu+m\tau_{\mathrm{rr}}(\dot a^\mu+a^2u^\mu/c^2)$，三个平行方向上的系数 $mc\dot\psi$、$F_\parallel$、$mc\tau_{\mathrm{rr}}\ddot\psi$ 相等即得上式。非线性的相对论运动学全部被快度吸收，而辐射反作用只作用于 $\dot\psi$。

令 $w(\tau)=\dot\psi(\tau)$，标量方程改写为
\begin{equation*}
 \dot w-{1\over\tau_{\mathrm{rr}}}w
 =-{F_\parallel(\tau)\over mc\tau_{\mathrm{rr}}}.
\end{equation*}
乘以积分因子 $e^{-\tau/\tau_{\mathrm{rr}}}$，得到
\begin{equation*}
 {\d\over\d\tau}
 \left(e^{-\tau/\tau_{\mathrm{rr}}}w\right)
 =-{e^{-\tau/\tau_{\mathrm{rr}}}\over mc\tau_{\mathrm{rr}}}
 F_\parallel(\tau).
\end{equation*}
从任意 $\tau_0$ 积分到 $\tau$，一般解为
\begin{equation*}
 w(\tau)
 =e^{(\tau-\tau_0)/\tau_{\mathrm{rr}}}
 \bigg[
 w(\tau_0)
 -{1\over mc\tau_{\mathrm{rr}}}
 \int_{\tau_0}^{\tau}
 e^{-(s-\tau_0)/\tau_{\mathrm{rr}}}F_\parallel(s)\d s
 \bigg].
\end{equation*}
方括号中的常数不受外力历史约束；若它非零，$w$ 含有 $e^{\tau/\tau_{\mathrm{rr}}}$ 的 runaway 模态。要求外力在未来关闭以后加速度不指数增长，相当于施加未来边界条件
\begin{equation*}
 \lim_{\tau\to+\infty}
 e^{-\tau/\tau_{\mathrm{rr}}}w(\tau)=0.
\end{equation*}
将积分因子关系式从 $\tau$ 积分到 $+\infty$：$\left[e^{-s/\tau_{\mathrm{rr}}}w(s)\right]_{\tau}^{\infty}=-\frac{1}{mc\tau_{\mathrm{rr}}}\int_{\tau}^{\infty}e^{-s/\tau_{\mathrm{rr}}}F_\parallel(s)\d s$，未来边界条件使上端为零，再令 $r=(s-\tau)/\tau_{\mathrm{rr}}$、$\d s=\tau_{\mathrm{rr}}\d r$，得到唯一的无 runaway 解
\begin{align*}
 \dot\psi(\tau)
 &={1\over mc\tau_{\mathrm{rr}}}
 \int_{\tau}^{\infty}
 e^{-(s-\tau)/\tau_{\mathrm{rr}}}F_\parallel(s)\d s\\
 &={1\over mc}
 \int_0^\infty e^{-r}
 F_\parallel(\tau+\tau_{\mathrm{rr}}r)\d r.
\end{align*}
同一个解的快度和动量为
\begin{equation*}
 \psi(\tau)=\psi(\tau_i)+\int_{\tau_i}^{\tau}\dot\psi(s)\d s,
 \qquad
 p(\tau)=mc\sinh\psi(\tau).
\end{equation*}
无 runaway 解的积分形式清楚显示：消去 runaway 的代价是当前加速度含有未来一个宽度约为 $\tau_{\mathrm{rr}}$ 的外力平均，这就是预加速的数学来源。

取最简单的阶跃力
\begin{equation*}
 F_\parallel(\tau)=F_0\Theta(\tau).
\end{equation*}
直接计算无 runaway 积分：对 $\tau<0$，阶跃在 $r=-\tau/\tau_{\mathrm{rr}}$ 处开启，$\int_{-\tau/\tau_{\mathrm{rr}}}^{\infty}e^{-r}\d r=e^{\tau/\tau_{\mathrm{rr}}}$；对 $\tau\geq0$ 全区间积分得 $1$。于是
\begin{equation*}
 \dot\psi(\tau)
 ={F_0\over mc}
 \begin{cases}
 e^{\tau/\tau_{\mathrm{rr}}},&\tau<0,\\
 1,&\tau\geq0.
 \end{cases}
\end{equation*}
若 $\psi(-\infty)=\psi_-$，再积分一次得到闭式
\begin{equation*}
 \psi(\tau)=\psi_-+{F_0\over mc}
 \begin{cases}
 \tau_{\mathrm{rr}}e^{\tau/\tau_{\mathrm{rr}}},&\tau<0,\\
 \tau+\tau_{\mathrm{rr}},&\tau\geq0,
 \end{cases}
 \qquad
 p(\tau)=mc\sinh\psi(\tau).
\end{equation*}
速度和快度在开关时刻连续，但加速度在外力开启以前已经按 $e^{\tau/\tau_{\mathrm{rr}}}$ 建立，这就是预加速。

为把“开启”和“关闭”放在同一个有限过程内，考虑矩形脉冲
\begin{equation*}
 F_\parallel(\tau)=F_0\bigl[\Theta(\tau)-\Theta(\tau-T)\bigr],
 \qquad T>0.
\end{equation*}
同一个无 runaway 积分给出：对 $0\leq\tau\leq T$ 只有前一个阶跃开启，$\int_0^{(T-\tau)/\tau_{\mathrm{rr}}}e^{-r}\d r=1-e^{-(T-\tau)/\tau_{\mathrm{rr}}}$；对 $\tau<0$ 积分区间为 $(-\tau/\tau_{\mathrm{rr}},(T-\tau)/\tau_{\mathrm{rr}})$，$\int e^{-r}\d r=e^{\tau/\tau_{\mathrm{rr}}}\left(1-e^{-T/\tau_{\mathrm{rr}}}\right)$；$\tau>T$ 时两个阶跃同时开启，积分为零。于是
\begin{equation*}
 \dot\psi(\tau)={F_0\over mc}
 \begin{cases}
 e^{\tau/\tau_{\mathrm{rr}}}
 \left(1-e^{-T/\tau_{\mathrm{rr}}}\right),&\tau<0,\\[2mm]
 1-e^{-(T-\tau)/\tau_{\mathrm{rr}}},&0\leq\tau\leq T,\\[2mm]
 0,&\tau>T.
 \end{cases}
\end{equation*}
逐段积分并要求 $\psi$ 连续，得到
\begin{equation*}
 \psi(\tau)=\psi_-+{F_0\over mc}
 \begin{cases}
 \tau_{\mathrm{rr}}e^{\tau/\tau_{\mathrm{rr}}}
 (1-e^{-T/\tau_{\mathrm{rr}}}),&\tau<0,\\[1mm]
 \tau+\tau_{\mathrm{rr}}
 \left[1-e^{-(T-\tau)/\tau_{\mathrm{rr}}}\right],&0\leq\tau\leq T,\\[1mm]
 T,&\tau>T.
 \end{cases}
\end{equation*}
粒子在脉冲到来前预加速，并在脉冲关闭前提前减小加速度；脉冲结束后没有 runaway，最终快度增量恰为 $F_0T/(mc)$。

若 $F_\parallel$ 在尺度 $T$ 上光滑变化且 $T\gg\tau_{\mathrm{rr}}$，可在无 runaway 积分中作 Taylor 展开：
\begin{align*}
 F_\parallel(\tau+\tau_{\mathrm{rr}}r)
 &=\sum_{n=0}^\infty
 {\tau_{\mathrm{rr}}^nr^n\over n!}
 {\d^nF_\parallel\over\d\tau^n},\\
 \dot\psi(\tau)
 &= {1\over mc}
 \sum_{n=0}^\infty
 {\tau_{\mathrm{rr}}^n\over n!}
 {\d^nF_\parallel\over\d\tau^n}
 \int_0^\infty e^{-r}r^n\d r\\
 &= {1\over mc}
 \sum_{n=0}^\infty
 \tau_{\mathrm{rr}}^n
 {\d^nF_\parallel\over\d\tau^n}.
\end{align*}
这里使用了 $\int_0^\infty e^{-r}r^n\d r=n!$。保留到一阶，
\begin{equation*}
 mc\dot\psi
 =F_\parallel+\tau_{\mathrm{rr}}\dot F_\parallel
 +O\!\left({\tau_{\mathrm{rr}}^2\over T^2}\right).
\end{equation*}
降阶式是精确无 runaway 解的低频渐近，而不是与 ALD 方程并列的另一条基本定律。对阶跃力，$\dot F_\parallel=F_0\delta(\tau)$，局域截断把平滑预加速层替换成一个 $\delta$ 冲量；这说明导数展开不能在不光滑外力处逐项使用。

上述降阶可直接写成协变形式。定义与四速度正交的投影算符
\begin{equation*}
 \Delta^\mu{}_{\nu}
 =\delta^\mu{}_{\nu}-{u^\mu u_\nu\over c^2}.
\end{equation*}
由 $u^2=c^2$ 求导得 $u\cdot a=0$，再对 $u\cdot a=0$ 求导得 $u\cdot\dot a=-a^2$，
\begin{align*}
 \Delta^\mu{}_{\nu}\dot a^\nu
 &=\dot a^\mu-{u^\mu\over c^2}(u\cdot\dot a)\\
 &=\dot a^\mu+{a^2\over c^2}u^\mu.
\end{align*}
因此 ALD 方程等价地写成
\begin{equation*}
 ma^\mu=f_{\mathrm{ext}}^\mu
 +m\tau_{\mathrm{rr}}\Delta^\mu{}_{\nu}\dot a^\nu.
\end{equation*}
在自力项中使用零阶关系 $ma^\mu=f_{\mathrm{ext}}^\mu+O(\tau_{\mathrm{rr}})$：对该关系沿世界线求导得 $\dot a^\nu=\frac{1}{m}\d f_{\mathrm{ext}}^\nu/\d\tau+O(\tau_{\mathrm{rr}})$，代入 $m\tau_{\mathrm{rr}}\Delta^\mu{}_\nu\dot a^\nu$，得到一阶降阶式
\begin{equation*}
 ma^\mu=f_{\mathrm{ext}}^\mu
 +\tau_{\mathrm{rr}}\Delta^\mu{}_{\nu}
 {\d f_{\mathrm{ext}}^\nu\over\d\tau}
 +O(\tau_{\mathrm{rr}}^2).
\end{equation*}
对 Lorentz 四力 $f_{\mathrm{ext}}^\mu=qF_{\mathrm{ext}}^\mu{}_{\nu}u^\nu$，先求导
\begin{equation*}
 {\d f_{\mathrm{ext}}^\mu\over\d\tau}
 =q(\partial_\alpha F_{\mathrm{ext}}^\mu{}_{\nu})u^\alpha u^\nu
 +qF_{\mathrm{ext}}^\mu{}_{\nu}a^\nu,
\end{equation*}
把该导数代入降阶式中的 $\tau_{\mathrm{rr}}\Delta^\mu{}_\nu\dot f_{\mathrm{ext}}^\nu$，第一项（$\partial_\alpha F$ 部分）保留，第二项（$F a$ 部分）用零阶关系 $a^\nu=(q/m)F_{\mathrm{ext}}^\nu{}_{\beta}u^\beta+O(\tau_{\mathrm{rr}})$ 替换，一般降阶式化为 Landau--Lifshitz 方程
\begin{align*}
 ma^\mu={}&qF_{\mathrm{ext}}^\mu{}_{\nu}u^\nu
 +q\tau_{\mathrm{rr}}
 \Delta^\mu{}_{\rho}
 (\partial_\alpha F_{\mathrm{ext}}^\rho{}_{\nu})u^\alpha u^\nu\\
 &+{q^2\tau_{\mathrm{rr}}\over m}
 \Delta^\mu{}_{\rho}
 F_{\mathrm{ext}}^\rho{}_{\nu}
 F_{\mathrm{ext}}^\nu{}_{\beta}u^\beta
 +O(\tau_{\mathrm{rr}}^2).
\end{align*}
投影算符使右端自动与 $u^\mu$ 正交。若用普通三力 $\vb F=\d\vb p/\d t$ 表示，并把自力项中的加速度替换为零阶加速度 $\vb a_0$，一般降阶式的空间部分可化为三维形式。记 $f^\mu=\gamma(\vb F\cdot\vb v/c,\vb F)$。因 $u_\mu f^\mu=0$，投影可改写成 $\Delta^{\mu}{}_{\nu}\dot f^\nu=\dot f^\mu+(a\cdot f/c^2)u^\mu$；由 $u^\mu=\gamma(c,\vb v)$ 与 $\d\gamma/\d t=\gamma^3(\vb v\cdot\vb a_0)/c^2$ 直接计算得 $a\cdot f=-\gamma^3\vb a_0\cdot\vb F$、$\d f^i/\d\tau=\gamma\dot\gamma F^i+\gamma^2\d F^i/\d t$。把 $ma^i=\gamma\d p^i/\d t$ 的空间分量除以 $\gamma$ 并代入 $\dot\gamma$，
\begin{align*}
 {\d p^i\over\d t}
 &=F^i+\tau_{\mathrm{rr}}\left[\gamma{\d F^i\over\d t}
 +{\gamma^3\over c^2}(\vb v\cdot\vb a_0)F^i
 -{\gamma^3\over c^2}v^i(\vb a_0\cdot\vb F)\right]\\
 &=F^i+\tau_{\mathrm{rr}}\left[\gamma{\d F^i\over\d t}
 -{\gamma^3\over c^2}\left(\vb a_0\times(\vb v\times\vb F)\right)^i\right],
\end{align*}
末行用到恒等式 $(\vb v\cdot\vb a_0)\vb F-\vb v(\vb a_0\cdot\vb F)=-\vb a_0\times(\vb v\times\vb F)$。写成向量形式：
\begin{equation*}
 {\d\vb p\over\d t}
 =\vb F
 +\tau_{\mathrm{rr}}
 \left[
 \gamma{\d\vb F\over\d t}
 -{\gamma^3\over c^2}
 \vb a_0\times(\vb v\times\vb F)
 \right]
 +O(\tau_{\mathrm{rr}}^2),
\end{equation*}
其中 $\d/\d t$ 是沿零阶轨道的全导数。这三个公式是同一降阶动力学的抽象投影形式、外电磁场形式和三维实验室形式。

恒定一维普通力 $F_\parallel=F_0$ 是一个有用检验。标量方程的无 runaway 解为 $\dot\psi=F_0/(mc)$，故 $\ddot\psi=0$，ALD 自力严格为零；与此同时 Lienard 功率非零。前面关于 Schott 项与远区辐射的守恒链式关系已经说明，这一辐射由 Schott 束缚动量的变化补偿。快度形式把“恒定固有加速度下有辐射而局域自力为零”的关系化成一行可检验的等式。这一节始终在点模型内部求解；下面保留源的有限形状，看同一个自作用理论如何给出附加惯性与延迟自力。

\subsubsection{有限尺寸正则化:附加惯性与延迟自力}

Dirac 计算说明，点粒子推迟场的共同奇异部分只通过质量重整化进入可观测运动方程，而推迟场与超前场的半差留下有限的 ALD 自力。现在反过来把点流 $J^\mu_{\mathrm{point}}$ 平滑为有限尺寸守恒电流。这样做不是另起一种理论，而是在同一个四维 Green 函数和同一个能动量守恒律中保留源的空间形状。附加质量、频率依赖惯性和延迟自力将作为这一平滑理论的低频、低速表示出现。

对任意光滑守恒电流，电磁能动量张量仍满足
\begin{equation*}
 T_{\mathrm{em}}^{\mu\nu}
 ={1\over\mu_0}
 \left(
 -F^{\mu\alpha}F^\nu{}_{\alpha}
 +{1\over4}g^{\mu\nu}F^{\alpha\beta}F_{\alpha\beta}
 \right),
 \qquad
 \partial_\mu T_{\mathrm{em}}^{\mu\nu}
 =-F^{\nu\lambda}J_\lambda.
\end{equation*}
这一关系由麦克斯韦方程直接推出：对张量逐项微分，并用 $\partial_\mu F^{\mu\alpha}=\mu_0J^\alpha$ 与 Bianchi 恒等式，
\begin{align*}
 \partial_\mu T_{\mathrm{em}}^{\mu\nu}
 &={1\over\mu_0}\left[-\partial_\mu F^{\mu\alpha}F^\nu{}_\alpha
 -F^{\mu\alpha}\partial_\mu F^\nu{}_\alpha
 +{1\over4}\partial^\nu(F^{\alpha\beta}F_{\alpha\beta})\right]\\
 &=-J^\alpha F^\nu{}_\alpha
 -{1\over2}F^{\mu\alpha}\left(\partial_\mu F^\nu{}_\alpha
 -\partial_\alpha F^\nu{}_\mu\right)
 +{1\over2}F^{\mu\alpha}\partial^\nu F_{\mu\alpha}\\
 &=-F^{\nu\lambda}J_\lambda,
\end{align*}
其中第二步第二项把导数换成反对称组合，利用了 $F^{\mu\alpha}$ 的反对称性；Bianchi 恒等式给出 $\partial_\mu F^\nu{}_\alpha-\partial_\alpha F^\nu{}_\mu=\partial^\nu F_{\mu\alpha}$，使末两项相消。
令 $T_{\mathrm{mat}}^{\mu\nu}$ 包含带电物质的惯性、维持其结构的约束应力及其他非电磁相互作用，并把总场写成
\begin{equation*}
 F^{\mu\nu}=F_{\mathrm{ext}}^{\mu\nu}+F_{\mathrm{self}}^{\mu\nu}.
\end{equation*}
物质与自场之间交换的四动量在总和中相消，因此局域守恒律可连续写成
\begin{align*}
 \partial_\mu T_{\mathrm{mat}}^{\mu\nu}
 &=F_{\mathrm{ext}}^{\nu\lambda}J_\lambda
   +F_{\mathrm{self}}^{\nu\lambda}J_\lambda,\\
 \partial_\mu
 \left(T_{\mathrm{mat}}^{\mu\nu}+T_{\mathrm{self}}^{\mu\nu}\right)
 &=F_{\mathrm{ext}}^{\nu\lambda}J_\lambda.
\end{align*}
第二行不是新的假设，而是第一行与自场能动量张量方程相加后的同一个守恒关系。取包围带电体的世界管 $\mathcal W$，以两张类空超曲面 $\Sigma_1,\Sigma_2$ 截断，把守恒关系在 $\Omega$ 上积分。Gauss 定理把 $\partial_\mu(T_{\mathrm{mat}}+T_{\mathrm{self}})^{\mu\nu}$ 的体积分化为 $\Sigma_2$、$-\Sigma_1$ 与管壁三部分的边界积分；物质场不穿过管壁，管壁项只剩自场部分。按 $P^\nu[\Sigma]$ 的定义整理，得到
\begin{align*}
 P_{\mathrm{mat}}^\nu[\Sigma_2]-P_{\mathrm{mat}}^\nu[\Sigma_1]
 +P_{\mathrm{self}}^\nu[\Sigma_2]-P_{\mathrm{self}}^\nu[\Sigma_1]
 +\int_{\partial\mathcal W}T_{\mathrm{self}}^{\mu\nu}\d\Sigma_\mu
 &=\int_\Omega F_{\mathrm{ext}}^{\nu\lambda}J_\lambda\d^4x,\\
 P^\nu[\Sigma]&={1\over c}\int_\Sigma T^{\mu\nu}\d\Sigma_\mu.
\end{align*}
世界管取得很大时，侧面通量是远区辐射四动量；世界管缩到源附近时，侧面和端面还含有随源一起运动的束缚场四动量。后面出现的附加惯性、Schott 动量和辐射反作用，都是同一自场能动量在不同尺度下的表示。

自场由推迟 Green 函数确定。在 Lorenz 规范下，
\begin{equation*}
 \Box D_{\mathrm{ret}}(x-x')=\delta^{(4)}(x-x'),
 \qquad
 A_{\mathrm{self}}^\mu(x)
 =\mu_0\int D_{\mathrm{ret}}(x-x')J^\mu(x')\d^4x',
\end{equation*}
故
\begin{equation*}
 F_{\mathrm{self}}^{\mu\nu}(x)
 =\mu_0\int
 \left[
 \partial^\mu D_{\mathrm{ret}}(x-x')J^\nu(x')
 -\partial^\nu D_{\mathrm{ret}}(x-x')J^\mu(x')
 \right]\d^4x'.
\end{equation*}
将自场代回 Lorentz 力密度，同一个自作用首先写成最一般的电流泛函
\begin{equation*}
 f_{\mathrm{self}}^\mu(x)
 =\mu_0J_\nu(x)\int
 \left[
 \partial^\mu D_{\mathrm{ret}}(x-x')J^\nu(x')
 -\partial^\nu D_{\mathrm{ret}}(x-x')J^\mu(x')
 \right]\d^4x'.
\end{equation*}
对于有限尺寸的光滑电流，该泛函有限且因果，但它依赖源的整个历史。若源可由若干集体坐标 $Q^A(t)$ 描述，在线性近似下将电流展开为
\begin{equation*}
 J^\mu(x;Q)=J_0^\mu(x)+J_A^\mu(x)Q^A+O(Q^2),
\end{equation*}
并把自力泛函投影到 $Q^A$。线性化后的自作用泛函是 $Q$ 的线性、平移不变泛函，因果性（$t$ 时刻的力只依赖 $t'<t$ 的历史）使这种泛函必为卷积形式，同一个自作用便改写为推迟记忆核
\begin{equation*}
 \mathcal F_A^{\mathrm{self}}(t)
 =\int_{-\infty}^tK_{AB}(t-t')Q^B(t')\d t',
 \qquad
 \widetilde{\mathcal F}_A^{\mathrm{self}}(\omega)
 =K_{AB}(\omega)\widetilde Q^B(\omega).
\end{equation*}
因 $K_{AB}(t<0)=0$，$K_{AB}(\omega)$ 在上半复频率平面解析，并满足
\begin{equation*}
 K_{AB}(-\omega)=K_{AB}(\omega)^*.
\end{equation*}
这两条性质分别表达因果性和实值性。

现取最重要的集体坐标，即整个电荷分布的平移 $\vb X(t)$。均匀位移不能产生自恢复力，均匀速度也不能产生耗散，因此平移核在 $\omega=0$ 至少有二阶零点。于是同一个平移自力可以依次写成
\begin{equation*}
 \widetilde F_{\mathrm{self},i}
 =K_{ij}(\omega)\widetilde X_j
 =\omega^2\mathcal M_{ij}^{\mathrm{em}}(\omega)\widetilde X_j
 =-\mathcal M_{ij}^{\mathrm{em}}(\omega)\widetilde a_j.
\end{equation*}
其中第二种与第三种写法由 $\vb a=\ddot{\vb X}$ 的频域关系 $\widetilde a_j=-\omega^2\widetilde X_j$ 等价。三种写法描述的是同一个物理量：第一种强调它是位移的线性响应，第二种抽出平移零模，第三种把剩余响应解释为频率依赖的附加惯性张量。一般形状的带电体只有张量 $\mathcal M_{ij}^{\mathrm{em}}$；只有施加旋转对称性以后才出现单一的“附加质量”。

下面由 Maxwell 方程直接计算惯性张量 $\mathcal M_{ij}^{\mathrm{em}}(\omega)$。设静止电荷密度为 $\rho_0(\vb x)$，取电荷中心为坐标原点，
\begin{equation*}
 q=\int\rho_0(\vb x)\d^3x,
 \qquad
 \int\vb x\rho_0(\vb x)\d^3x=0.
\end{equation*}
作小振幅刚性平移并保留 $\vb X$ 的一次项，
\begin{equation*}
 \rho(\vb x,t)
 =\rho_0(\vb x)-\vb X(t)\cdot\nabla\rho_0(\vb x),
 \qquad
 \vb J(\vb x,t)=\rho_0(\vb x)\dot{\vb X}(t).
\end{equation*}
采用 Fourier 约定
\begin{equation*}
 g(\vb x,t)=\int{\d^3k\d\omega\over(2\pi)^4}
 \widetilde g(\vb k,\omega)e^{i\vb k\cdot\vb x-i\omega t},
 \qquad
 \widetilde\rho_0(\vb k)=\int\rho_0(\vb x)e^{-i\vb k\cdot\vb x}\d^3x,
\end{equation*}
在该约定下 $\nabla\mapsto i\vb k$、$\partial_t\mapsto -i\omega$；对 $\rho$ 的表达式作变换并对 $\rho_0$ 分部积分（$\widetilde{\nabla\rho_0}=i\vb k\widetilde\rho_0$），连续性方程化为
\begin{equation*}
 \widetilde{\delta\rho}
 =-i\widetilde\rho_0\,\vb k\cdot\widetilde{\vb X},
 \qquad
 \widetilde{\vb J}
 =-i\omega\widetilde\rho_0\widetilde{\vb X}.
\end{equation*}
波动方程 $\Box\Phi=-\rho/\varepsilon_0$、$\Box\vb A=-\mu_0\vb J$ 的 Fourier 形式为 $(-k^2+\omega^2/c^2)\widetilde\Phi=-\widetilde\rho/\varepsilon_0$ 等，解出推迟标势和矢势（分母中的 $-i0$ 是推迟 Green 函数 $1/(k^2-(\omega/c)^2-i0)$ 的因果处方）
\begin{equation*}
 \widetilde\Phi
 ={\widetilde{\delta\rho}\over
 \varepsilon_0[k^2-(\omega/c)^2-i0]},
 \qquad
 \widetilde{\vb A}
 ={\mu_0\widetilde{\vb J}\over
 k^2-(\omega/c)^2-i0},
\end{equation*}
由 $\vb E=-\nabla\Phi-\partial_t\vb A$，频域中 $\widetilde{\delta E}_i=-ik_i\widetilde{\delta\Phi}+i\omega\widetilde A_i$，代入上式并利用 $\mu_0=1/(\varepsilon_0c^2)$，运动引起的电场扰动为
\begin{equation*}
 \widetilde{\delta E}_i
 ={\widetilde\rho_0\over\varepsilon_0[k^2-(\omega/c)^2-i0]}
 \left({\omega^2\over c^2}\delta_{ij}-k_i k_j\right)
 \widetilde X_j.
\end{equation*}
总自力的一次项不是只积分 $\rho_0\delta\vb E$：位移密度 $\delta\rho$ 还作用于静态自场 $\vb E_0$，故
\begin{equation*}
 \widetilde F_{\mathrm{self},i}
 =\int\left(\rho_0\widetilde{\delta E}_i
 +\widetilde{\delta\rho}\,E_{0i}\right)\d^3x.
\end{equation*}
第二项与第一项的纵向部分精确相消：由 $\widetilde E_{0i}=-ik_i\widetilde\rho_0/(\varepsilon_0k^2)$，第二项按 Parseval 恒等式给出 $\int\frac{|\widetilde\rho_0|^2k_ik_j\widetilde X_j}{\varepsilon_0k^2}\frac{\d^3k}{(2\pi)^3}$；而 $\frac{\omega^2}{c^2}\delta_{ij}-k_ik_j=\frac{\omega^2}{c^2}\left(\delta_{ij}-\frac{k_ik_j}{k^2}\right)-k_ik_j\left(1-\frac{\omega^2}{c^2k^2}\right)$ 的纵向部分满足 $-\frac{k_ik_j(1-\omega^2/c^2k^2)}{k^2-\omega^2/c^2}=-\frac{k_ik_j}{k^2}$，恰好与之抵消。这正是整块电荷刚性平移不产生净自力的平移对称性。引入横向投影
\begin{equation*}
 P_{ij}^{\mathrm T}(\vb k)=\delta_{ij}-{k_i k_j\over k^2},
\end{equation*}
利用 Parseval 恒等式整理后，同一个自力依次化为
\begin{align*}
 \widetilde F_{\mathrm{self},i}
 &= {\omega^2\widetilde X_j\over\varepsilon_0c^2}
 \int{\d^3k\over(2\pi)^3}
 {\left|\widetilde\rho_0(\vb k)\right|^2P_{ij}^{\mathrm T}(\vb k)
 \over k^2-(\omega/c)^2-i0}\\
 &=-\mathcal M_{ij}^{\mathrm{em}}(\omega)\widetilde a_j,
\end{align*}
其中
\begin{equation*}
 \mathcal M_{ij}^{\mathrm{em}}(\omega)
 ={1\over\varepsilon_0c^2}
 \int{\d^3k\over(2\pi)^3}
 {\left|\widetilde\rho_0(\vb k)\right|^2P_{ij}^{\mathrm T}(\vb k)
 \over k^2-(\omega/c)^2-i0}.
\end{equation*}
这个惯性张量是线性平移问题中自力泛函的闭式化结果。它仍保留任意电荷形状、有限光行时和完整的线性辐射反作用，尚未使用长波展开。

把附加惯性分成实部和虚部，
\begin{equation*}
 \mathcal M_{ij}^{\mathrm{em}}
 =\mathcal M_{ij}'+i\mathcal M_{ij}'',
\end{equation*}
实部改变同相加速度响应，虚部产生不可逆功率。由
\begin{equation*}
 {1\over x-i0}=\mathcal P{1\over x}+i\pi\delta(x),
\end{equation*}
同一个 $\mathcal M_{ij}^{\mathrm{em}}$ 的虚部先写成三维壳层积分；$\delta(k^2-k_0^2)=\delta(k-k_0)/(2k_0)$ 使径向积分 $\int_0^\infty k^2\d k\,\delta(k^2-k_0^2)=k_0/2$（$k_0=\omega/c$），随后完成径向积分，得到更透明的角积分形式
\begin{align*}
 \mathcal M_{ij}''(\omega)
 &= {\pi\over\varepsilon_0c^2}
 \int{\d^3k\over(2\pi)^3}
 |\widetilde\rho_0(\vb k)|^2P_{ij}^{\mathrm T}(\vb k)
 \delta\!\left(k^2-{\omega^2\over c^2}\right)\\
 &= {\omega\over16\pi^2\varepsilon_0c^3}
 \int\d\Omega_{\vb n}
 \left|\widetilde\rho_0\!\left({\omega\over c}\vb n\right)\right|^2
 \left(\delta_{ij}-n_i n_j\right),
 \qquad \omega>0.
\end{align*}
对简谐加速度 $\vb a(t)=\Re(\vb a_0e^{-i\omega t})$，平均 $\langle F_iv_i\rangle=\frac12\Re(\widetilde F_i^*\widetilde v_i)$；由 $\widetilde v_i=i\widetilde a_i/\omega$、$\widetilde F_i=-\mathcal M_{ij}\widetilde a_j$ 与 $\mathcal M=\mathcal M'+i\mathcal M''$，实部 $\mathcal M'$ 的贡献为纯虚，自场所做平均功为
\begin{equation*}
 \left\langle\vb F_{\mathrm{self}}\cdot\vb v\right\rangle
 =-{1\over2\omega}a_{0i}^*\mathcal M_{ij}''(\omega)a_{0j}.
\end{equation*}
因此同一个虚部也可由辐射功率定义：
\begin{align*}
 \left\langle P_{\mathrm{rad}}\right\rangle
 &={1\over2\omega}a_{0i}^*\mathcal M_{ij}''(\omega)a_{0j}\\
 &={1\over32\pi^2\varepsilon_0c^3}
 \int\d\Omega_{\vb n}
 \left|\widetilde\rho_0\!\left({\omega\over c}\vb n\right)\right|^2
 \left|\vb n\times\vb a_0\right|^2.
\end{align*}
第一行把功率写成响应函数的耗散部分，第二行把同一功率写成前面频谱分析中的远区角积分，其中用了 $a_{0i}^*(\delta_{ij}-n_in_j)a_{0j}=|\vb n\times\vb a_0|^2$。两种计算相等，说明附加惯性和辐射谱不是两个独立问题。

因惯性张量是推迟响应，$K(\omega)$ 在上半复频率平面解析并在无穷远衰减。对 $K(\omega')/(\omega'-\omega)$ 沿实轴加无穷远半圆作 Cauchy 积分，得 $K(\omega)=\frac{1}{i\pi}\mathcal P\int_{-\infty}^{\infty}\frac{K(\omega')}{\omega'-\omega}\d\omega'$；用实值性 $K(-\omega)=K^*(\omega)$ 折到正频率，虚实部满足色散关系
\begin{equation*}
 \mathcal M_{ij}'(\omega)
 =\mathcal M_{ij}'(\infty)
 +{2\over\pi}\mathcal P\int_0^\infty
 {\omega'\mathcal M_{ij}''(\omega')\over\omega'^2-\omega^2}\d\omega'.
\end{equation*}
对有限尺寸源，$\mathcal M_{ij}'(\infty)=0$（谱表示中积分以 $k\sim|\omega|/c$ 为特征尺标，形状因子在高频衰减）；且在 $\omega=0$ 处 $\mathcal M_{ij}''(0)=0$，故 $\mathcal M_{ij}(0)=\mathcal M_{ij}'(0)$。令 $\omega=0$，同一个静态附加惯性又化为整个辐射谱密度的和规则
\begin{equation*}
 \mathcal M_{ij}^{\mathrm{em}}(0)
 ={2\over\pi}\int_0^\infty
 {\mathcal M_{ij}''(\omega')\over\omega'}\d\omega'.
\end{equation*}
因此，零频附加质量并不是只由近场“局域地”决定；它也等于全部频率辐射耗散的色散积分。

惯性张量在零频处为
\begin{equation*}
 \mathcal M_{ij}^{\mathrm{em}}(0)
 ={1\over\varepsilon_0c^2}
 \int{\d^3k\over(2\pi)^3}
 {|\widetilde\rho_0(\vb k)|^2\over k^2}
 \left(\delta_{ij}-{k_i k_j\over k^2}\right).
\end{equation*}
对指标取迹。横向投影的迹为 $\operatorname{tr}P^{\mathrm T}=3-1=2$；又由 $U_{\mathrm{es}}=\frac12\int\rho_0\Phi_0\d^3x$ 与 $\widetilde\Phi_0=\widetilde\rho_0/(\varepsilon_0k^2)$ 经 Parseval 恒等式，同一个积分可识别成静电能
\begin{equation*}
 U_{\mathrm{es}}
 ={1\over2\varepsilon_0}
 \int{\d^3k\over(2\pi)^3}{|\widetilde\rho_0(\vb k)|^2\over k^2},
\end{equation*}
便有连续的等式链
\begin{equation*}
 \operatorname{tr}\boldsymbol{\mathcal M}^{\mathrm{em}}(0)
 ={2\over\varepsilon_0c^2}
 \int{\d^3k\over(2\pi)^3}{|\widetilde\rho_0|^2\over k^2}
 ={4U_{\mathrm{es}}\over c^2}.
\end{equation*}
这是一条适用于任意形状电荷分布的迹定理。若再取球对称，旋转不变性使
\begin{equation*}
 \mathcal M_{ij}^{\mathrm{em}}(0)=\delta m_{\mathrm{dyn}}\delta_{ij},
\end{equation*}
于是迹定理的左端成为 $3\delta m_{\mathrm{dyn}}$，同一个结论才进一步简化为
\begin{equation*}
 \delta m_{\mathrm{dyn}}
 ={1\over3}\operatorname{tr}\boldsymbol{\mathcal M}^{\mathrm{em}}(0)
 ={4U_{\mathrm{es}}\over3c^2}.
\end{equation*}
这说明 $4/3$ 系数首先是一般附加惯性张量在球对称下的标量化结果，而不是可以在任意源模型中预先假定的常数。这是 $4/3$ 的第一次出现。

接下来在球对称条件下把同一个频率响应进一步化简，并展示 Jackson--Abraham--Lorentz 推导中实空间展开、闭式求和和形状因子表示之间的关系。设电荷分布作非相对论刚性平移，瞬时静止系中 $\vb J(\vb x,t)=\rho(\vb x)\vb v(t)$。定义电磁动量对机械方程的贡献为
\begin{equation*}
 {\d\vb p_{\mathrm{em}}\over\d t}
 \equiv-\vb F_{\mathrm{self}}.
\end{equation*}
将推迟势中的源量在 $t-R/c$ 处展开，
\begin{equation*}
 [g]_{\mathrm{ret}}
 =\sum_{n=0}^\infty{(-1)^n\over n!}
 \left({R\over c}\right)^n g^{(n)}(t),
 \qquad R=|\vb x-\vb x'|,
\end{equation*}
代入自场 Lorentz 力。推迟展开的 $n$ 阶项中，标势的 $\vb R/R^2$ 部分经连续性方程 $\dot\rho=-\nabla'\cdot\vb J$ 并入矢势项，$n$ 阶积分中出现
\begin{equation*}
 \vb J(\vb x')-{\vb R\over n+2}\nabla'\cdot\vb J(\vb x').
\end{equation*}
对第二项分部积分（$\partial'_jR_i=-\delta_{ij}$、$\partial'_jR=-R_j/R$），在双重积分内它等价于
\begin{equation*}
 {n+1\over n+2}\vb J
 -{n-1\over n+2}{(\vb J\cdot\vb R)\vb R\over R^2}.
\end{equation*}
再代入 $\vb J=\rho\vb v$。球对称使所有 $\vb R$ 方向等权：$\langle R_iR_j\rangle_\Omega=\alpha\delta_{ij}$，两侧缩并 $\langle R^2\rangle_\Omega=3\alpha$ 定出 $\alpha=R^2/3$，故
\begin{equation*}
 \left\langle{(\vb v\cdot\vb R)\vb R\over R^2}\right\rangle_{\Omega_R}
 ={1\over3}\vb v.
\end{equation*}
于是上面的冗长向量组合在同一个积分中化成
\begin{equation*}
 {n+1\over n+2}\vb J
 -{n-1\over n+2}{(\vb J\cdot\vb R)\vb R\over R^2}
 \longrightarrow {2\over3}\rho(\vb x')\vb v.
\end{equation*}
自场动量率因此得到只含标量距离矩的展开
\begin{equation*}
 {\d\vb p_{\mathrm{em}}\over\d t}
 ={1\over6\pi\varepsilon_0c^2}
 \sum_{n=0}^\infty{(-1)^n\over n!c^n}
 \vb v^{(n+1)}(t)
 \iint\rho(\vb x)\rho(\vb x')R^{n-1}\d^3x\d^3x'.
\end{equation*}
级数是同一个球对称自作用的时间域导数展开。其 $n=0$ 项是附加惯性，$n=1$ 项是形状无关的辐射反作用：
\begin{align*}
 {\d\vb p_{\mathrm{em}}\over\d t}
 &=\left[{1\over6\pi\varepsilon_0c^2}
 \iint{\rho(\vb x)\rho(\vb x')\over R}\d^3x\d^3x'\right]\vb a
 -{q^2\over6\pi\varepsilon_0c^3}\dot{\vb a}+\cdots\\
 &={4U_{\mathrm{es}}\over3c^2}\vb a
 -{q^2\over6\pi\varepsilon_0c^3}\dot{\vb a}+\cdots.
\end{align*}
其中 $n=0$ 项系数用 $U_{\mathrm{es}}=\frac12\int\rho\Phi\d^3x=\frac{1}{8\pi\varepsilon_0}\iint\frac{\rho(\vb x)\rho(\vb x')}{|\vb x-\vb x'|}\d^3x\d^3x'$ 化为 $4U_{\mathrm{es}}/(3c^2)$，$n=1$ 项系数用 $\iint\rho(\vb x)\rho(\vb x')\d^3x\d^3x'=q^2$ 给出 $q^2/(6\pi\varepsilon_0c^3)$。由于 $\vb F_{\mathrm{self}}=-\d\vb p_{\mathrm{em}}/\d t$，上式立即给出熟悉的惯性反力和 Abraham--Lorentz 项，但它们此时只是完整级数的前两项。

对级数作 Fourier 变换，$\vb v^{(n+1)}\mapsto(-i\omega)^{n+1}\widetilde{\vb v}$。系数部分求和（$(-1)^n(-i)^n=i^n$）：
\begin{equation*}
 \sum_{n=0}^\infty{(-1)^n\over n!c^n}(-i\omega)^{n+1}R^{n-1}
 ={-i\omega\over R}\sum_{n=0}^\infty{(i\omega R/c)^n\over n!}
 ={-i\omega\over R}e^{i\omega R/c},
\end{equation*}
幂级数中的距离矩重新求和为
\begin{equation*}
 \sum_{n=0}^\infty{1\over n!}
 \left({i\omega R\over c}\right)^n
 ={e^{i\omega R/c}},
\end{equation*}
因此同一个无穷导数级数精确重写成
\begin{align*}
 \widetilde{\dot{\vb p}}_{\mathrm{em}}
 &=-i\omega M_{\mathrm{em}}(\omega)\widetilde{\vb v},\\
 M_{\mathrm{em}}(\omega)
 &={1\over6\pi\varepsilon_0c^2}
 \iint\rho(\vb x)\rho(\vb x')
 {e^{i\omega R/c}\over R}\d^3x\d^3x'.
\end{align*}
级数形式和闭式形式是同一个物理量的两种写法；前者适合长波展开，后者保留全部光行时。

再定义球对称形状因子
\begin{equation*}
 \widetilde\rho(\vb k)=qf(k),
 \qquad f(0)=1.
\end{equation*}
用外向波 Green 函数的标准展开（$(\nabla^2+k_0^2)G=\delta^{(3)}$ 的 Fourier 解）
\begin{equation*}
 {e^{ik_0R}\over R}
 =4\pi\int{\d^3k\over(2\pi)^3}{e^{i\vb k\cdot\vb R}\over k^2-k_0^2-i0},
 \qquad k_0={\omega\over c},
\end{equation*}
代入闭式并对 $d^3x\,d^3x'$ 作卷积，$|\widetilde\rho(\vb k)|^2=q^2|f(k)|^2$，角向积分 $\int\d\Omega=4\pi$ 后，同一个 $M_{\mathrm{em}}(\omega)$ 化成谱表示
\begin{equation*}
 M_{\mathrm{em}}(\omega)
 ={q^2\over3\pi^2\varepsilon_0c^2}
 \int_0^\infty
 {k^2|f(k)|^2\over k^2-(\omega/c)^2-i0}\d k.
\end{equation*}
实空间表示和谱表示分别把同一附加惯性写成电荷对之间的推迟相互作用和空间频率的叠加：前者显示有限传播时延，后者显示源形状如何截断高频响应。

若裸机械质量为 $m_0$，机械动量和电磁动量相加后，外力方程为
\begin{equation*}
 -i\omega\mathsf M(\omega)\widetilde{\vb v}
 =\widetilde{\vb F}_{\mathrm{ext}},
 \qquad
 \mathsf M(\omega)=m_0+M_{\mathrm{em}}(\omega).
\end{equation*}
物理低频质量是同一个有效质量在零频处的值，
\begin{equation*}
 m\equiv\mathsf M(0)
 =m_0+{4U_{\mathrm{es}}\over3c^2}.
\end{equation*}
用低频质量消去裸质量后，外力方程可改写成更清楚、紫外收敛性更好的形式
\begin{align*}
 \mathsf M(\omega)
 &=m+{1\over6\pi\varepsilon_0c^2}
 \iint\rho\rho'
 {e^{i\omega R/c}-1\over R}\d^3x\d^3x'\\
 &=m+{q^2\omega^2\over3\pi^2\varepsilon_0c^4}
 \int_0^\infty
 {|f(k)|^2\over k^2-(\omega/c)^2-i0}\d k.
\end{align*}
第二行由第一行代入谱表示并相减：$\frac{e^{i\omega R/c}-1}{R}$ 对应频域差 $\frac{1}{k^2-k_0^2-i0}-\frac{1}{k^2}=\frac{k_0^2}{k^2(k^2-k_0^2-i0)}$。两行仍是同一个 $\mathsf M(\omega)$：减去零频值以后，原来线性发散的静态质量被隔离，剩余的频率依赖部分多了两个高频幂次的收敛性。

作为与 Dirac 点粒子结果的交叉检验，现在在已经减去零频质量的响应中取点形状 $f(k)=1$。利用
\begin{equation*}
 \int_0^\infty{\d k\over k^2-(\omega/c)^2-i0}
 ={i\pi c\over2\omega},
 \qquad \omega>0,
\end{equation*}
该值是主值 $\mathcal P\int_0^\infty\frac{\d k}{k^2-k_0^2}=0$（被积函数关于 $k_0$ 奇对称）与 $k=k_0$ 处留数 $\pi i/(2k_0)$ 之和。同一个有效质量化为
\begin{equation*}
 \mathsf M_{\mathrm{point}}(\omega)
 =m+i{q^2\omega\over6\pi\varepsilon_0c^3}
 =m(1+i\omega\tau_{\mathrm{rr}}),
 \qquad
 \tau_{\mathrm{rr}}={q^2\over6\pi\varepsilon_0mc^3}.
\end{equation*}
代回外力方程并逆变换，得到
\begin{equation*}
 m\vb a
 =\vb F_{\mathrm{ext}}+m\tau_{\mathrm{rr}}\dot{\vb a}.
\end{equation*}
因此这里重新得到的 Abraham--Lorentz 方程与 Dirac 正则场的低速极限完全相同；它也是同一个有限尺寸有效质量 $\mathsf M(\omega)$ 在先减去静态质量、再取点形状以后得到的特殊形式。

点形状抹掉了高频结构。为看清有限尺寸如何恢复时延，取半径为 $a$ 的均匀薄球壳，
\begin{equation*}
 f(k)={\sin ka\over ka},
 \qquad
 U_{\mathrm{es}}={q^2\over8\pi\varepsilon_0a},
 \qquad
 \delta m={4U_{\mathrm{es}}\over3c^2}
 ={q^2\over6\pi\varepsilon_0ac^2}.
\end{equation*}
将这个 $f(k)$ 代入形状因子公式。利用下列标准积分（第一式由 $\sin^2ka=\frac12(1-\cos2ka)$ 与 $\mathcal P\int_0^\infty\frac{\cos(bk)}{k^2-k_0^2}\d k=-\frac{\pi}{2k_0}\sin(bk_0)$ 得到；第二式取 $\Im\frac{1}{k^2-k_0^2-i0}=\pi\delta(k^2-k_0^2)$ 的留数）
\begin{equation*}
 \mathcal P\int_0^\infty{\sin^2(ka)\over k^2-k_0^2}\d k
 ={\pi\over4k_0}\sin(2k_0a),
 \qquad
 \Im\int_0^\infty{\sin^2(ka)\over k^2-k_0^2-i0}\d k
 ={\pi\over2k_0}\sin^2(k_0a),
\end{equation*}
得到
\begin{equation*}
 M_{\mathrm{em}}^{\mathrm{shell}}(\omega)
 =\delta m\,{e^{i\zeta_a}-1\over i\zeta_a},
 \qquad
 \zeta_a={2\omega a\over c}.
\end{equation*}
核对：$\frac{e^{i\zeta}-1}{i\zeta}=\frac{\sin\zeta}{\zeta}+i\frac{1-\cos\zeta}{\zeta}$，代入 $\delta m=q^2/(6\pi\varepsilon_0ac^2)$ 后实部、虚部与上面两个积分一致；零频极限 $M_{\mathrm{em}}^{\mathrm{shell}}(0)=\delta m$（$\int_0^\infty\frac{\sin^2x}{x^2}\d x=\pi/2$）。同一个球壳自力先在频域写成
\begin{align*}
 \widetilde{\vb F}_{\mathrm{self}}
 &=-M_{\mathrm{em}}^{\mathrm{shell}}(\omega)\widetilde{\vb a}\\
 &={q^2\over12\pi\varepsilon_0ca^2}
 \left(e^{2i\omega a/c}-1\right)\widetilde{\vb v},
\end{align*}
逆变换后仍是同一个自力，但其因果结构变得直接可见：
\begin{equation*}
 \vb F_{\mathrm{self}}(t)
 ={q^2\over12\pi\varepsilon_0ca^2}
 \left[\vb v\!\left(t-{2a\over c}\right)-\vb v(t)\right].
\end{equation*}
最后在 $|\omega|a/c\ll1$ 下展开这个精确延迟力。令 $\tau=2a/c$，Taylor 展开 $\vb v(t-\tau)-\vb v(t)=-\tau\vb a+\frac{\tau^2}{2}\dot{\vb a}-\frac{\tau^3}{6}\ddot{\vb a}+\cdots$，得
\begin{align*}
 \vb F_{\mathrm{self}}
 &={q^2\over12\pi\varepsilon_0ca^2}
 \left[-{2a\over c}\vb a
 +{1\over2}\left({2a\over c}\right)^2\dot{\vb a}
 -{1\over6}\left({2a\over c}\right)^3\ddot{\vb a}+\cdots\right]\\
 &=-\delta m\,\vb a
 +{q^2\over6\pi\varepsilon_0c^3}\dot{\vb a}
 -{q^2a\over9\pi\varepsilon_0c^4}\ddot{\vb a}+\cdots.
\end{align*}
频率响应、频域力、时域延迟和局域渐近是同一个球壳自作用的四种写法。只有最后一行出现三阶导数，因此 runaway 问题属于局域截断，而不是精确延迟核的必然性质。

外力为零时频域方程 $-i\omega[m_0+M_{\mathrm{em}}^{\mathrm{shell}}(\omega)]\widetilde{\vb v}=0$ 给出无外力模态：
\begin{equation*}
 \omega\left[m_0+M_{\mathrm{em}}^{\mathrm{shell}}(\omega)\right]=0,
\end{equation*}
除平移零模（$\omega=0$）外，稳定性要求响应函数在上半频率平面无零点。在上半平面取 $\omega=i\xi$（$\xi>0$），此时 $e^{i\zeta_a}=e^{-2\xi a/c}$，$m_0+M_{\mathrm{em}}^{\mathrm{shell}}(i\xi)=m_0+\delta m\frac{1-e^{-2\xi a/c}}{2\xi a/c}$ 是 $\xi$ 的单调减函数（$\frac{1-e^{-x}}{x}$ 在 $x>0$ 递减），从 $m$（$\xi\to0$）降到 $m_0$（$\xi\to\infty$）；除非 $m_0\geq0$，否则它必然穿越零点，给出 $e^{\xi t}$ 增长的不稳定模态。因此无零点的条件可写成
\begin{equation*}
 m_0\geq0
 \qquad\Longleftrightarrow\qquad
 a\geq c\tau_{\mathrm{rr}},
\end{equation*}
其中 $m=m_0+\delta m$。若在保持物理质量 $m$ 固定时把 $a$ 压到 $c\tau_{\mathrm{rr}}$ 以下，就必须让裸质量变成负值，精确响应的稳定性随之丧失。点粒子三阶方程中的异常解由这一极限结构产生，而不是由“辐射必然导致反因果”产生。

光滑电荷给出同一母式的另一种形象特例。取
\begin{equation*}
 \rho_0(r)={q\over\pi^{3/2}a^3}e^{-r^2/a^2},
 \qquad
 f(k)=e^{-a^2k^2/4},
 \qquad
 x={\omega a\over\sqrt2c}.
\end{equation*}
将它代入形状因子公式：常数部分 $\int_0^\infty e^{-a^2k^2/2}\d k=\frac1a\sqrt{\frac\pi2}$，频率依赖部分用
\begin{equation*}
 \mathcal P\int_0^\infty{e^{-\alpha k^2}\over k^2-b^2}\d k
 =-{\pi\over2b}e^{-\alpha b^2}\operatorname{erfi}(b\sqrt\alpha),
\end{equation*}
同一个 $M_{\mathrm{em}}(\omega)$ 化为闭式
\begin{equation*}
 M_{\mathrm{em}}^{\mathrm G}(\omega)
 ={q^2\over3\pi^2\varepsilon_0c^2}
 \left[
 {\sqrt\pi\over\sqrt2a}
 -{\pi\omega\over2c}e^{-x^2}\operatorname{erfi}(x)
 +i{\pi\omega\over2c}e^{-x^2}
 \right].
\end{equation*}
其零频值仍是 $4U_{\mathrm{es}}/(3c^2)$，而虚部被 $e^{-\omega^2a^2/(2c^2)}$ 平滑截断。球壳和 Gaussian 分布具有不同高频核，却共享同一个低频辐射系数，这正是一般角积分公式的内容。

到这里，$4U_{\mathrm{es}}/(3c^2)$ 已经由自力响应得到。现在从场动量重新计算同一个零频惯性，便可看清传统 $4/3$ 问题的来源。将静电构型以小速度 $\vb v$ 平移，最低阶磁场为
\begin{equation*}
 \vb B={1\over c^2}\vb v\times\vb E.
\end{equation*}
其中 $(\vb E\times\vb B)_i=\frac{1}{c^2}\epsilon_{ijk}\epsilon_{klm}E_jv_lE_m=\frac{1}{c^2}(E^2\delta_{ij}-E_iE_j)v_j$（用到 $\epsilon_{ijk}\epsilon_{klm}=\delta_{il}\delta_{jm}-\delta_{im}\delta_{jl}$），电磁动量为
\begin{align*}
 P_{\mathrm{em},i}
 &=\varepsilon_0\int(\vb E\times\vb B)_i\d^3x\\
 &={\varepsilon_0\over c^2}
 \int\left(E^2\delta_{ij}-E_iE_j\right)\d^3x\,v_j.
\end{align*}
静电场的 Fourier 分量为（$\widetilde\Phi_0=\widetilde\rho_0/(\varepsilon_0k^2)$，$\vb E_0=-\nabla\Phi_0$）
\begin{equation*}
 \widetilde E_i(\vb k)
 =-{ik_i\over\varepsilon_0k^2}\widetilde\rho_0(\vb k).
\end{equation*}
将它代入电磁动量并用 Parseval 恒等式（$\int E_iE_j\d^3x=\int\frac{k_ik_j|\widetilde\rho_0|^2}{\varepsilon_0^2k^4}\frac{\d^3k}{(2\pi)^3}$，$\int E^2\d^3x=\int\frac{|\widetilde\rho_0|^2}{\varepsilon_0^2k^2}\frac{\d^3k}{(2\pi)^3}$），方括号中的惯性张量化成
\begin{align*}
 {\varepsilon_0\over c^2}
 \int\left(E^2\delta_{ij}-E_iE_j\right)\d^3x
 &={1\over\varepsilon_0c^2}
 \int{\d^3k\over(2\pi)^3}
 {|\widetilde\rho_0|^2\over k^2}
 \left(\delta_{ij}-{k_i k_j\over k^2}\right)\\
 &=\mathcal M_{ij}^{\mathrm{em}}(0).
\end{align*}
因此场动量和前面由自力得到的零频附加惯性是同一个量：
\begin{equation*}
 P_{\mathrm{em},i}
 =\mathcal M_{ij}^{\mathrm{em}}(0)v_j.
\end{equation*}
球对称时，上式立即给出
\begin{equation*}
 \vb P_{\mathrm{em}}
 ={4U_{\mathrm{es}}\over3c^2}\vb v.
\end{equation*}
这不是第二个无关的 $4/3$ 结论，而是同一 $\mathcal M^{\mathrm{em}}(0)$ 分别在自力和场动量中的出现，这是 $4/3$ 的第二次出现。

然而纯电磁场不是孤立守恒系统。对任意静止构型，在小速度 boost 后，同一实验室同时面上的总动量为
\begin{equation*}
 P^i={v_j\over c^2}
 \left[
 \delta_{ij}\int T^{00}\d^3x
 +\int T^{ij}\d^3x
 \right]+O(v^3/c^3).
\end{equation*}
其来源：小速度 boost 下 $T^{i0}=T'^{i0}+\frac{v_j}{c}T'^{ij}+\frac{v^i}{c}T'^{00}+O(v^2/c^2)$，静止构型 $T'^{i0}=0$，由 $P^i=\frac1c\int T^{i0}\d^3x$ 即得该式。它表明，能量密度给出 $E/c^2$，积分应力还会给出额外惯性。纯电磁球对称场满足
\begin{equation*}
 \int T_{\mathrm{em}}^{ij}\d^3x
 ={U_{\mathrm{es}}\over3}\delta_{ij},
\end{equation*}
所以总系数成为 $4U_{\mathrm{es}}/(3c^2)$，这是 $4/3$ 的第三次出现。稳定带电体还必须含有抵抗 Coulomb 排斥的非电磁应力 $T_{\mathrm P}^{\mu\nu}$。对局域、静态且总能动量守恒的系统，
\begin{equation*}
 \partial_\mu T_{\mathrm{tot}}^{\mu\nu}=0,
 \qquad
 T_{\mathrm{tot}}^{\mu\nu}=T_{\mathrm{em}}^{\mu\nu}+T_{\mathrm P}^{\mu\nu},
\end{equation*}
由
\begin{equation*}
 \partial_j(x^iT_{\mathrm{tot}}^{jk})
 =T_{\mathrm{tot}}^{ik}+x^i\partial_jT_{\mathrm{tot}}^{jk}
\end{equation*}
在全空间积分。静态平衡 $\partial_jT^{jk}=0$ 使右端第二项为零；左端是散度的全空间积分，化为无穷远面积分，对有限源消失，得到 Laue 条件
\begin{equation*}
 \int T_{\mathrm{tot}}^{ij}\d^3x=0.
\end{equation*}
把 Laue 条件代回 boost 公式，总动量才化为
\begin{equation*}
 \vb P_{\mathrm{tot}}
 ={E_{\mathrm{tot}}\over c^2}\vb v.
\end{equation*}
因此 $U_{\mathrm{es}}/c^2$ 与 $4U_{\mathrm{es}}/(3c^2)$ 并非互相矛盾的两个“质量”。后者是非守恒电磁子系统在实验室同时面上的动力学惯性，前者是能量项；只有把维持稳定的物质应力一并加入，上式才给出协变的总质量。

薄球壳提供这一应力闭合的熟悉图像。由 Gauss 定律 $E(a^+)=q/(4\pi\varepsilon_0a^2)$，壳外电场对球面产生向外 Maxwell 压力
\begin{equation*}
 p_{\mathrm{em}}
 ={\varepsilon_0E^2(a^+)\over2}
 ={q^2\over32\pi^2\varepsilon_0a^4}.
\end{equation*}
若以表面张力 $\mathcal T$ 维持平衡，Laplace 关系 $2\mathcal T/a=p_{\mathrm{em}}$ 给出
\begin{equation*}
 \mathcal T={q^2\over64\pi^2\varepsilon_0a^3}.
\end{equation*}
该张力的积分应力抵消电磁场的积分应力，使 Laue 条件成立；至于稳定物质本身含有多少静能，则仍由具体微观模型决定。至此有限尺寸正则化的图像闭合：自力的 $4/3$、场动量的 $4/3$ 与 boost 公式中的 $4/3$ 是同一个 $\mathcal M^{\mathrm{em}}(0)$ 的不同出场方式，而有效质量 $\mathsf M(\omega)$ 把附加惯性、频率依赖惯性与延迟自力统一在同一响应函数中。下一节把这个 $\mathsf M(\omega)$ 用于辐射振子，看它如何同时决定附加质量、频移、线宽与散射截面。
